\begin{center}
\LARGE
A Topology-Independent Model for Railway Interlocking
                             Systems
\end{center}

\begin{center}
\Large
Eugenio Roanes-Lozano, Luis M. Laita

\end{center}

\noindent
Date:  July 18th (Thursday)\\
Time:  14:25-14:50\\

\begin{center}
\large
Abstract
\end{center}

Two decision models for interlocking systems are presented. The first for the case that the turnouts
have spring switches and the second for the case that the turnouts should not be trailed through a
switch set against. In both of them, the algorithm is independent from the topology of the station
and is based on the calculation of accessibility in an oriented graph. Surprisingly brief
implementations in Maple V of both models have been developed. The package analyzes any
change in the position of the switches and the clearances given by the semaphores and supervises if
it is safe to authorize the change. In the second model, which turnouts could be trailed through
switches set against is also studied. To finish with, an example is given. 

